% META: {"id": "TNSI-ALGO-CO-008", "chapitre": "TNSI-ALGORITHMIQUE", "type_objet": "corrige", "exercice_ref": "TNSI-ALGO-EX-008", "capacites": ["T-ALGO-03"], "capacites_codes": ["C11"], "mode_creation": "ex_nihilo", "sources_inspiration": [], "status": "generated"}
% BEGIN-VERIFY
% altitudes = [120, 180, 260, 410, 620, 900, 780, 540, 300, 150]
% assert all(altitudes[i] < altitudes[i + 1] for i in range(5))
% assert all(altitudes[i] > altitudes[i + 1] for i in range(5, 9))
% appels = []
% def sommet(altitudes, gauche, droite):
%     appels.append((gauche, droite))
%     if gauche == droite:
%         return gauche
%     milieu = (gauche + droite) // 2
%     if altitudes[milieu] < altitudes[milieu + 1]:
%         return sommet(altitudes, milieu + 1, droite)
%     return sommet(altitudes, gauche, milieu)
% assert sommet(altitudes, 0, 9) == 5
% assert altitudes[5] == 900 == max(altitudes)
% assert appels == [(0, 9), (5, 9), (5, 7), (5, 6), (5, 5)]
% assert len(appels) == 5 and 9 > 5
% assert altitudes[4] < altitudes[5] and altitudes[5] > altitudes[6]
% appels.clear()
% deux = [10, 90, 20, 80, 15]
% trouve = sommet(deux, 0, 4)
% assert deux[trouve] == 80 and max(deux) == 90
% assert deux[trouve - 1] < deux[trouve] and deux[trouve + 1] < deux[trouve]
% assert 2 ** 4 > len(altitudes)
% END-VERIFY
\begin{corrige}{TNSI-ALGO-EX-008}
\textbf{1.} Les cinq premiers écarts sont positifs
($120<180<260<410<620<900$) et les quatre suivants négatifs
($900>780>540>300>150$) : le relevé est bien strictement croissant puis
strictement décroissant. Un parcours complet demanderait $9$ comparaisons, une
par élément après le premier.

\textbf{2.} La fonction ne conserve qu'une moitié à chaque appel :
\begin{python}
def sommet(altitudes, gauche, droite):
    if gauche == droite:
        return gauche
    milieu = (gauche + droite) // 2
    if altitudes[milieu] < altitudes[milieu + 1]:
        return sommet(altitudes, milieu + 1, droite)
    return sommet(altitudes, gauche, milieu)
\end{python}

\textbf{3.} Si \lstinline{altitudes[m] < altitudes[m+1]}, on est encore dans la
phase de montée à l'indice $m$. Or le relevé monte strictement \emph{jusqu'au}
sommet : tant qu'on monte, tout ce qui précède est plus bas que ce qui suit.
Aucun indice inférieur ou égal à $m$ ne peut donc porter le maximum, et on peut
écarter toute la première moitié sans l'examiner.

C'est cette propriété — décider avec une seule comparaison quelle moitié jeter —
qui rend la méthode applicable ici.

\textbf{4.} L'appel \lstinline{sommet(altitudes, 0, 9)} renvoie l'indice $5$,
d'altitude $900$~m. Il a fallu $5$ appels, correspondant aux intervalles
$(0\,;\,9)$, $(5\,;\,9)$, $(5\,;\,7)$, $(5\,;\,6)$ et $(5\,;\,5)$ : l'intervalle
est divisé par deux à chaque étape. Contre $9$ comparaisons pour le parcours
complet, l'écart est modeste sur dix points, mais il croît comme $n$ contre
$\log_2 n$ : sur un relevé d'un million de points, ce serait un million contre
une vingtaine.

\textbf{5.} La différence est essentielle. Pour le maximum d'une liste
\emph{quelconque}, diviser en deux moitiés oblige à explorer les \textbf{deux} :
le maximum peut être n'importe où, et on ne peut rien écarter. Le nombre total de
comparaisons reste de l'ordre de $n$ ; la récursivité ne fait que réorganiser le
travail.

Ici, l'hypothèse sur la forme du relevé permet d'\emph{éliminer} une moitié.
Diviser pour régner ne fait gagner que lorsqu'une propriété des données permet de
ne pas descendre dans les deux branches — c'est ce qui sépare la dichotomie d'un
simple découpage récursif.

\textbf{6.} Sur \lstinline{[10, 90, 20, 80, 15]}, la fonction renvoie l'indice
$3$, d'altitude $80$, alors que le maximum de la liste vaut $90$.

Elle n'est pas fautive pour autant : la valeur trouvée est bien supérieure à ses
deux voisines, c'est un \emph{maximum local}. C'est tout ce que la comparaison de
la question~3 peut garantir dès lors que le relevé n'a plus une seule montée
suivie d'une seule descente.

La conclusion est que l'hypothèse de départ n'est pas un ornement de l'énoncé :
elle fait partie de la spécification de la fonction. Employée hors de son domaine
de validité, celle-ci rend un résultat plausible et faux — sans lever la moindre
erreur.
\end{corrige}
